\section{Reproducibilidad y contribuciones}

\subsection{Comando, entradas y ejecución} \label{sec:repro}

El código está en el repositorio público del grupo (https://github.com/jorge-mardones-data/electricity-demand-forecasting.git), versión \pend{commit de la entrega final}. La única entrada es el paquete del caso (\texttt{data/raw/demanda\_horaria.csv}). Con XeLaTeX instalado (no se instala con \texttt{pip}), desde la raíz:
\begin{verbatim}
python -m pip install -r requirements.txt
python analisis.py
\end{verbatim}
\texttt{analisis.py} ejecuta auditoría, base analítica, validación (con los hiperparámetros fijados), sensibilidades, evaluación en 2019, bootstrap, diagnóstico, tablas, figuras y la compilación de este informe; \texttt{--busqueda-completa} repite la búsqueda de hiperparámetros (varias horas). Registra entorno y tiempos en \texttt{results/run\_info.json} y la diferencia con los resultados previos en \texttt{results/concordancia.csv}.

\textbf{Ejecuciones.} Los resultados provienen de una ejecución completa el 9 de noviembre de 2026 en un solo computador (Computadora A): Windows 11, Intel Core i5-14600K (20 núcleos lógicos), 31.7 GB de RAM, Python 3.11.9, NumPy 2.4.6, pandas 2.3.3, scikit-learn 1.7.2 (Pedregosa et al., 2011) y matplotlib 3.10.7, con semillas fijas (\texttt{random\_state} $=42$). La ejecución tardó 12.2 minutos, y la búsqueda de hiperparámetros del Random Forest, fuera de ella, 178 minutos. El 10 de noviembre se repitió en la computadora de otro integrante (Computadora B: AMD Ryzen 9600X, 12 núcleos lógicos, Ubuntu 24.04 en WSL2): en CPU, los pasos 3 a 6 tardaron 10.6 minutos frente a 11.9 en A. En una GPU NVIDIA RTX 5060 de 8 GB no se aceleraron; fueron algo más lentos (anexo, tabla \ref{tab:tiempos_equipos}).

\textbf{Concordancia y copia limpia.} Frente a los resultados previos del repositorio (generados con los notebooks del grupo), las diferencias son menores a 1,4\,\% de la escala de cada archivo: solo cambia visiblemente el Random Forest cíclico de 2019 (NMAE macro de 0,02493 a 0,02490), plausiblemente por cálculos de punto flotante (seno y coseno) distintos entre equipos o versiones de librerías, lo que no se verificó. El 9 de noviembre de 2026 se copió la carpeta (sin \texttt{.git} ni registros previos) a otra ubicación y se ejecutó \texttt{python analisis.py} en el mismo computador: terminó sin errores en 13,8 minutos con la misma concordancia. En B terminó sin errores y con diferencias menores a 1,4\,\%.

\subsection{Contribuciones del equipo}

\textbf{Jorge Mardones:} creación y estructuración del repositorio y todo el contenido de su primera versión (preparación y auditoría de los datos, base analítica, protocolo experimental, diseño de la validación y desarrollo de los modelos, notebooks 02 y 03). \textbf{Benjamín Ruiz:} secciones 3.3, 4.1 y 4.2 del enunciado (sensibilidades, evaluación en la reserva y métricas; notebooks 04 y 05 y \texttt{src/modeling.py}) y la ejecución de la comparación de modelos. \textbf{Matías Pineda:} redacción del informe e integración de todos los resultados, análisis complementarios (notebook 06), tablas y figuras (\texttt{src/make\_report\_assets.py}), \texttt{analisis.py}, \texttt{README.md} y \texttt{requirements.txt}. \textbf{Diego Paredes:} presentación del trabajo (\texttt{reports/presentacion/main.tex}, \texttt{style.tex} y \texttt{src/make\_presentation\_assets.py}, con sus figuras y tablas), revisión y corrección de detalles estructurales del informe, y verificar resultados.

\subsection{Declaración de uso de inteligencia artificial}

\textbf{Herramienta y fechas:} Claude Code (Anthropic), 8 y 9 de noviembre de 2026. \textbf{Propósito y partes afectadas:} borrador y recorte de este informe, integración de los resultados de 2019, bootstrap, diagnóstico y sensibilidades, \texttt{src/make\_report\_assets.py}, \texttt{src/run\_validation.py}, \texttt{analisis.py}, \texttt{README.md}, \texttt{requirements.txt}, el notebook 06, una celda de exportación en el notebook 03 y dos notas corregidas del notebook 05. \textbf{Comprobaciones:} tablas y figuras se generan desde \texttt{results/}; el bootstrap recalculado reproduce el del notebook 05; las cifras del texto se contrastaron con la ejecución única; \texttt{analisis.py} se ejecutó completo dos veces; los nombres de los feriados provienen del calendario oficial. Los borradores contenían errores de interpretación que la revisión detectó (sección 3.3).

\textbf{Herramienta y fechas:} ChatGPT (OpenAI), 4-8 de noviembre de 2026. \textbf{Propósito y partes afectadas:} herramienta de apoyo para estructura del repositorio inicial, exploración y auditoria de datos, construcción de la base analítica, desarrollo computacional en evaluacion y comparación de los cinco modelos iniciales. Se usó códigos principalmente en \texttt{docs/protocolo\_experimental.md}, en \texttt{src/audit\_data.py}, \texttt{src/build\_dataset.py} y en \texttt{notebooks/02\_experimental\_design.ipynb} y \texttt{notebooks/03\_model\_development.ipynb}. Los códigos y cualquier tipo de sugerencia fué revisada y adaptada por Jorge Mardones, responsable de las decisiones metodológicas iniciales del proyecto.

\textbf{Herramienta y fechas:} Claude Code (Anthropic), 8-10 de noviembre de 2026. \textbf{Propósito y partes afectadas:} Diseño, cohesión y estructura de presentación (\texttt{reports/presentacion/main.tex} y \texttt{style.tex}) a partir de los resultados obtenidos, y creación de \texttt{src/make\_presentation\_assets.py} para generar figuras y tablas con formatos aptos para exposición en sala. Revisión y corrección de la estructura de flotantes que desplazaban secciones del anexo técnico, y generación de la tabla de sensibilidad por país agregada a \texttt{make\_report\_assets.py}, para luego incluirla en el anexo (\texttt{content/anexo.tex}, \texttt{style.cls})

\textbf{Herramienta y fechas:} Claude (Anthropic, modo Cowork), 8--10 de noviembre de 2026: primera versión de los notebooks 04 y 05 y de \texttt{src/modeling.py} (verificado: reproduce el NMAE del notebook 03) y medición de tiempos CPU/GPU.
